\documentclass[UTF8,a4paper,12pt]{ctexart}

\usepackage[left=2.35cm,right=2.35cm,top=2.35cm,bottom=2.35cm]{geometry}
\usepackage{graphicx}
\usepackage{amsmath,amssymb}
\usepackage{booktabs}
\usepackage{tabularx}
\usepackage{longtable}
\usepackage{array}
\usepackage[table]{xcolor}
\usepackage{caption}
\usepackage{subcaption}
\usepackage{float}
\usepackage{enumitem}
\usepackage{listings}
\usepackage{tcolorbox}
\usepackage{fancyhdr}
\usepackage{etoolbox}
\usepackage{tocloft}
\usepackage[hidelinks]{hyperref}

\graphicspath{{figures/}}
\definecolor{deepblue}{HTML}{1E4E8C}
\definecolor{softblue}{HTML}{EAF3FF}
\definecolor{accent}{HTML}{D15B35}
\definecolor{goodgreen}{HTML}{1F8A5B}
\definecolor{lightgray}{HTML}{F5F7FA}
\definecolor{codebg}{HTML}{F7F9FC}
\definecolor{codeframe}{HTML}{D7E3F3}
\hypersetup{
  colorlinks=true,
  linkcolor=deepblue,
  urlcolor=deepblue,
  citecolor=deepblue,
  pdfauthor={负责人、控制设计成员、仿真验证成员},
  pdftitle={工业控制系统设计与仿真分析报告}
}

\pagestyle{fancy}
\fancyhf{}
\fancyhead[L]{\textcolor{deepblue}{工业控制系统设计与仿真分析报告}}
\fancyfoot[C]{\thepage}
\renewcommand{\headrulewidth}{0.4pt}
\setlength{\headheight}{15pt}

\setlist[itemize]{leftmargin=2em,itemsep=0.2em,topsep=0.2em}
\setlist[enumerate]{leftmargin=2em,itemsep=0.2em,topsep=0.2em}
\renewcommand{\arraystretch}{1.25}
\renewcommand{\thesection}{\chinese{section}}
\renewcommand{\thesubsection}{\arabic{section}.\arabic{subsection}}
\ctexset{
  section = {
    name = {,、},
    format = \zihao{3}\bfseries,
    beforeskip = 0pt,
    afterskip = 1.0em
  },
  subsection = {
    format = \zihao{4}\bfseries,
    beforeskip = 0.9em,
    afterskip = 0.45em
  }
}
\setlength{\parindent}{2em}
\setlength{\parskip}{0.25em}
\pretocmd{\section}{\ifnum\value{section}>0\clearpage\fi}{}{}

\lstdefinestyle{reportcode}{
  basicstyle=\ttfamily\small,
  backgroundcolor=\color{codebg},
  frame=single,
  rulecolor=\color{codeframe},
  breaklines=true,
  columns=fullflexible,
  keepspaces=true,
  keywordstyle=\color{deepblue}\bfseries,
  commentstyle=\color{goodgreen},
  stringstyle=\color{accent},
  numbers=left,
  numberstyle=\tiny\color{gray},
  xleftmargin=1.5em,
  framexleftmargin=1.2em
}
\lstset{style=reportcode}
\newcommand{\figref}[1]{图~\ref{#1}}
\newcommand{\tabref}[1]{表~\ref{#1}}
\renewcommand{\contentsname}{目\quad 录}
\setcounter{tocdepth}{2}
\setlength{\cftbeforesecskip}{0.8em}
\setlength{\cftsecindent}{0em}
\setlength{\cftsecnumwidth}{2.2em}
\setlength{\cftsubsecindent}{2.2em}
\setlength{\cftsubsecnumwidth}{3.2em}
\renewcommand{\cftsecfont}{\bfseries}
\renewcommand{\cftsecpagefont}{\bfseries}

\begin{document}
\begin{titlepage}
  \thispagestyle{empty}
  \centering
  \vspace*{2.4cm}
  {\zihao{1}\bfseries\textcolor{deepblue}{工业控制系统设计与仿真分析报告}\par}
  \vspace{1.2cm}
  {\zihao{-3} 控制器设计、性能验证与工程评价\par}
  \vfill
  \begin{tabular}{p{3.1cm}p{8.5cm}}
    项目/课程：& Co-Sight 工业控制系统设计智能体\\[0.65em]
    团队：& 示例团队\\[0.65em]
    负责人：& 负责人\\[0.65em]
    成员：& 负责人、控制设计成员、仿真验证成员\\[0.65em]
    班级/部门：& 工业控制系统设计组\\[0.65em]
    提交日期：& 2026-08-03\\[0.65em]
  \end{tabular}
  \vspace*{1cm}
\end{titlepage}
\clearpage
\pagenumbering{Roman}
\pagestyle{plain}
\tableofcontents
\clearpage
\pagenumbering{arabic}
\pagestyle{fancy}
\section*{摘要}
本文以一个具有单位负反馈结构的工业控制对象为例，建立对象模型，分析原系统的稳态精度与动态性能，比较候选控制方案，并通过统一仿真指标验证最终方案。报告将控制公式、仿真配置、性能数据和证据来源绑定到同一份可复现的结构化输入中，最终生成 HTML 预览和 LaTeX 源码。
\par
\noindent\textbf{关键词：} 工业控制系统；控制器设计；性能指标；仿真验证；工程评价
\clearpage
\section{题目背景与被控对象}\label{sec:model}
先明确系统边界、输入输出关系和评价目标，再进入控制器设计。
\par
被控对象采用单位负反馈结构。设计阶段只使用离线模型和历史数据，现场控制器保持隔离。所有结果都应标明模型版本、采样周期、输入信号和仿真时长。


\begin{equation}
G(s)=\frac{250}{s(s+2)(s+40)(s+45)}
\label{eq:plant}
\end{equation}

\begin{equation}
T(s)=\frac{G_c(s)G(s)}{1+G_c(s)G(s)}
\label{eq:closed_loop}
\end{equation}

\begin{tcolorbox}[title=\textbf{验证主线},colback=softblue,colframe=deepblue]
先判断原系统是否满足稳态指标，再比较单纯提高增益与引入相位补偿的代价，最后用同一套指标确认推荐方案。
\end{tcolorbox}

\section{性能指标与方案约束}\label{sec:requirements}
本示例将超调量、上升时间、2\% 调节时间和速度误差常数作为主要评价指标。工程约束同时包括控制器参数范围、采样周期、执行器饱和边界和离线仿真安全边界。


\begin{table}[H]
  \centering
  \caption{设计指标与验收阈值}
  \label{tab:requirements}
  \begin{tabularx}{\textwidth}{*{4}{>{\raggedright\arraybackslash}X}}
    \toprule
    指标 & 符号 & 单位 & 验收条件 \\
    \midrule
    超调量 & Mp & \% & 小于 20 \\
    上升时间 & tr & s & 小于 0.5 \\
    调节时间 & ts & s & 小于 1.2 \\
    速度误差常数 & Kv & 无量纲 & 大于等于 10 \\
    \bottomrule
  \end{tabularx}
\end{table}

\begin{enumerate}
  \item 统一模型、输入和仿真时长后比较所有候选方案。
  \item 记录每个指标的计算方法和容差设置。
  \item 只有通过硬约束的候选方案才能进入工程推荐。
\end{enumerate}

\section{控制方案设计与机理}\label{sec:design}
仅增加比例增益可以提高低频增益，但可能降低相位裕度并放大超调。超前校正通过在目标交越频率附近提供正相位，改善快速性与稳定裕度之间的折中。


\begin{equation}
G_c(s)=1483.7\frac{s+3.5}{s+33.75}
\label{eq:controller}
\end{equation}

\begin{tcolorbox}[title=\textbf{安全边界},colback=softblue,colframe=goodgreen]
控制器参数只写入离线仿真配置，渲染工具不具备 PLC、DCS、机器人或现场执行器访问能力。
\end{tcolorbox}

\section{仿真验证与性能对比}\label{sec:validation}
仿真采用统一阶跃输入、统一仿真时长和 2\% 调节带宽。结果表同时保留基准、候选和推荐方案，便于审查设计决策而不是只展示最终数字。


\begin{table}[H]
  \centering
  \caption{不同控制方案性能指标对比}
  \label{tab:performance}
  \begin{tabularx}{\textwidth}{*{6}{>{\raggedright\arraybackslash}X}}
    \toprule
    方案 & Kv & 上升时间/s & 调节时间/s & 超调量/\% & 是否满足 \\
    \midrule
    原系统 & 0.0694 & 30.406 & 54.691 & 0.000 & 否 \\
    比例增益 K=144 & 10.000 & 0.260 & 7.368 & 68.061 & 否 \\
    超前校正 & 10.685 & 0.167 & 0.829 & 19.067 & 是 \\
    \bottomrule
  \end{tabularx}
\end{table}

推荐方案同时满足四项硬指标，但超调量距离上限较近。因此工程部署前仍应进行参数摄动、采样周期和执行器饱和条件下的鲁棒性验证。


\section{工程评价与结论}\label{sec:conclusion}
结果表明，单纯提高增益虽然能够改善稳态精度，但会显著牺牲动态品质；超前校正通过频域补偿获得了更好的综合折中。最终推荐方案应连同模型版本、参数配置、仿真脚本和证据引用一起归档。


\begin{tcolorbox}[title=\textbf{推荐结论},colback=softblue,colframe=goodgreen]
推荐将超前校正方案作为离线设计候选，并在进入任何现场实施流程前完成独立复核和安全审批。
\end{tcolorbox}

\begin{lstlisting}[language=Python,caption={性能指标复核示例}]
metrics = {'overshoot_pct': 19.067, 'rise_time_s': 0.167, 'settling_time_s': 0.829, 'Kv': 10.685}\npassed = metrics['overshoot_pct'] < 20 and metrics['rise_time_s'] < 0.5 and metrics['settling_time_s'] < 1.2 and metrics['Kv'] >= 10\nprint(passed)
\end{lstlisting}

\clearpage
\section*{参考文献}
\begin{thebibliography}{99}
\bibitem{ref_control_textbook} 课程教材编写组. 自动控制原理课程教材. 课程资料库. 2026. 频域校正与性能指标章节
\end{thebibliography}
\appendix
\section{可复现性说明}
正式报告应在此处补充仿真脚本路径、模型版本、运行环境、图表文件哈希和原始结果文件路径。


\end{document}
